\section{Conclusiones y trabajos futuros}
\begin{spacing}{1.5}

\subsection{Conclusiones}
\label{sec:conclusiones}

A partir de los resultados obtenidos se extraen las siguientes conclusiones, organizadas según el grado de cumplimiento de los objetivos específicos planteados.

\subsubsection{Sobre los objetivos alcanzados}

\begin{itemize}
    \item \textbf{Viabilidad del modelado paramétrico (objetivo 1).} La superposición de cinco funciones Gaussianas asimétricas, complementada con un término de línea base, reproduce la morfología del ciclo P-QRS-T en ritmo sinusal normal con una correlación de Pearson media de 0,9986 y un RMSE medio de 0,0496 sobre 2.238 latidos. Las desviaciones estándar respectivas, 0,0020 y 0,0162, indican que el desempeño es estable a lo largo del registro y no depende de latidos favorables aislados.

    \item \textbf{Eficacia de la optimización con restricciones (objetivo 2).} El algoritmo L-BFGS-B \cite{byrd1995limited} con restricciones de caja permitió ajustar los veintidós parámetros del modelo a cada latido individual sin generar morfologías biológicamente imposibles. El resultado depende, no obstante, de que las restricciones estén expresadas en las mismas unidades que la señal a la que se aplican. La primera implementación derivó sus cotas de amplitud de rangos fisiológicos en milivoltios y las aplicó a latidos normalizados por puntuación Z: el 99,8\% de los latidos convergió exactamente sobre la cota superior de la amplitud de la onda R y, en promedio, 11,8 de los 20 parámetros quedaron determinados por un límite en lugar de por los datos. Expresadas las cotas de amplitud como fracción del pico de cada latido, ese promedio desciende a 1,8 de 22.

    \item \textbf{Valor diagnóstico del residuo frente a la métrica agregada.} La correlación de Pearson de la primera iteración, 0,9647, era compatible con un modelo cuya mayoría de parámetros no estaba siendo ajustada por los datos. Fueron dos indicadores auxiliares —el conteo de parámetros saturados en sus cotas y el sesgo medio del residuo— los que localizaron ambos defectos y guiaron una corrección que redujo el RMSE en un 87,6\%. Reportar métricas agregadas de fidelidad sin instrumentar el residuo habría dejado ese margen sin detectar.

    \item \textbf{Interpretabilidad frente a los modelos generativos profundos.} Cada parámetro del modelo posee un significado fisiológico directo, propiedad de la que carecen las arquitecturas basadas en redes generativas antagónicas \cite{brophy2023generative}. Esta característica es una ventaja estructural del enfoque, si bien su valor práctico para la aumentación de datos requiere validación adicional.
\end{itemize}

\subsubsection{Sobre los objetivos no alcanzados}

La honestidad metodológica exige declarar que tres de los cinco objetivos específicos no fueron completados en esta etapa:

\begin{itemize}
    \item El motor de reglas fisiológicas basado en las tablas de Rijnbeek \cite{rijnbeek2001normal} y de la AHA \cite{rautaharju2009aha} (objetivo 3) no fue implementado. En consecuencia, el generador no produce, por ahora, variabilidad demográfica por edad ni por sexo.

    \item El modelado de alteraciones patológicas (objetivo 4), específicamente la elevación del segmento ST y la Fibrilación Auricular, tampoco fue abordado.

    \item La validación (objetivo 5) se completó parcialmente: se calcularon la correlación de Pearson y el RMSE, pero no se realizó el análisis en el dominio de la frecuencia comprometido en la propuesta.
\end{itemize}

\subsubsection{Sobre la hipótesis}

La hipótesis planteada sostenía que es posible construir un generador de señales ECG sintéticas basado en Gaussianas y optimización numérica que reproduzca la morfología de señales reales con alta fidelidad estructural, permitiendo además la parametrización controlada por edad, sexo y condición patológica.

La evidencia obtenida valida la primera parte de la hipótesis y deja la segunda sin verificar. La reproducción morfológica quedó demostrada para ritmo sinusal normal en población adulta. La parametrización demográfica y patológica no fue puesta a prueba, por lo que no puede afirmarse ni refutarse con los datos disponibles.

Corresponde asimismo precisar la expresión ``alta fidelidad''. La primera iteración del modelo alcanzaba una correlación de 0,9647 con un RMSE de 0,3985, combinación que describe una reconstrucción morfológicamente correcta pero con una discrepancia de amplitud apreciable: la correlación de Pearson es invariante ante transformaciones afines y, por sí sola, sobreestima el desempeño del generador. En el modelo corregido ambas métricas son consistentes entre sí —0,9986 y 0,0496— de modo que la expresión resulta defendible, si bien únicamente para ritmo sinusal normal en población adulta.

\subsection{Trabajos Futuros}

Las líneas de trabajo inmediatas se desprenden directamente de las limitaciones identificadas en la sección de resultados:

\begin{enumerate}
    \item \textbf{Redefinir la ventana de extracción.} La ventana de 128 muestras a 360 Hz cubre 356 ms, mientras que un complejo P-QRS-T normal se extiende entre 600 y 800 ms. Las ondas P y T quedan parcialmente fuera del recorte y sus anchos permanecen determinados por sus cotas en cerca de la mitad de los latidos, pese a la corrección de las restricciones. Ampliar la ventana obliga a redefinir la unidad de análisis y a recalcular las métricas reportadas, por lo que constituye una decisión metodológica y no un ajuste de implementación.

    \item \textbf{Extender la validación al conjunto de la base.} Los resultados presentados corresponden a un único registro de los 48 que componen MIT-BIH. La generalización del método exige repetir el procedimiento sobre el conjunto completo y reportar la variabilidad entre registros.

    \item \textbf{Implementar el motor de reglas demográficas.} Integrar las tablas de Rijnbeek \cite{rijnbeek2001normal} y de la AHA \cite{rautaharju2009aha} para modular los intervalos PR, QRS y QT según edad y sexo, y validar las señales resultantes contra rangos de referencia clínicos.

    \item \textbf{Modelar alteraciones patológicas.} Incorporar la elevación del segmento ST mediante funciones de lesión y la Fibrilación Auricular mediante estocasticidad del intervalo R-R con anulación de la onda P.

    \item \textbf{Completar la validación en el dominio de la frecuencia.} Contrastar los espectros de potencia de las señales sintéticas y reales, según lo comprometido en el quinto objetivo específico.

    \item \textbf{Validar la utilidad para aumentación de datos.} La justificación última de este trabajo es mejorar el entrenamiento de clasificadores de arritmias. Corresponde entrenar un clasificador con datos reales, repetir el entrenamiento incorporando las señales sintéticas y medir la diferencia de desempeño. Sin este experimento, la utilidad práctica del generador permanece como hipótesis.
\end{enumerate}

\end{spacing}
